% Source: https://artofproblemsolving.com/community/c3h2594653_trig_questions_that_i_cant_understand_the_solutions_to_
% Stable identity: aops-topic-2594653-q01
Let $u=\dfrac{x+y}2$ and $v=\dfrac{x-y}2$, so that $x=u+v$ and $y=u-v$. By the angle addition formulas,
\[\sin(u+v)=\sin u\cos v+\cos u\sin v,\qquad\sin(u-v)=\sin u\cos v-\cos u\sin v.\]
Adding,
\[\sin x+\sin y=\sin(u+v)+\sin(u-v)=2\sin u\cos v=2\sin\dfrac{x+y}2\cos\dfrac{x-y}2.\ \blacksquare\]

Geometric picture: let $P=(\cos x,\sin x)$ and $Q=(\cos y,\sin y)$ on the unit circle, and $R=P+Q$. Then $OPRQ$ is a rhombus, so $R$ lies on the angle bisector of $\angle POQ$, at angle $\dfrac{x+y}2$, and $OR=2\,OM=2\cos\dfrac{x-y}2$, where $M$ is the midpoint of chord $PQ$. The height of $R$ is $\sin x+\sin y$, which is the identity.

\begin{asy}
usepackage("amsmath");
size(4.5cm);
defaultpen(fontsize(10pt));
real x=75,y=20;pair O=(0,0),P=dir(x),Q=dir(y),Mm=(P+Q)/2,R=P+Q;
draw(arc(O,1,-10,100),gray);draw((-0.2,0)--(1.9,0),Arrow(4));draw((0,-0.2)--(0,1.8),Arrow(4));
draw(arc(O,0.2,0,x));label("$x$",0.28*dir(x/2+8),fontsize(7pt));draw(arc(O,0.45,0,y));label("$y$",0.55*dir(y/2),fontsize(7pt));
draw(O--P,linewidth(0.8));draw(O--Q,linewidth(0.8));draw(P--R--Q,dashed);draw(O--R,dotted);draw(P--Q,gray);
draw(R--(R.x,0),dotted);
dot(P);dot(Q);dot(Mm);dot(R);dot(O);
label("$O$",O,SW,fontsize(8pt));label("$P$",P,NW,fontsize(8pt));label("$Q$",Q,SE,fontsize(8pt));label("$M$",Mm,SE,fontsize(8pt));label("$R$",R,NE,fontsize(8pt));
\end{asy}
